%% kapa: exam=2 round=28 year=2017 period=3 subject=compensation_law subject_ko=감정평가 및 보상법규 seq=2 number=2 points=30
%% source: Q-Net 감정평가사 기출문제 (cst00309) artlSeq=5201712 file=raw/1008/5201712/2017년 감정평가사 2차시험 기출문제(감정평가 및 보상법규).pdf
%% conversion: kordoc 4.24.51 -> tools/gen2.py (auto); review=proofread-sonnet; flags=none
\begin{longquestion}{【문제 2】}
도지사 A는 “X국가산업단지 내 국도 대체우회도로개설사업”(이하 ‘이 사건 개발사업’이라 함)의 실시계획을 승인·고시하고, 사업시행자로 B시의 시장을 지정하였다. B시의 시장은 이 사건 개발사업을 시행함에 있어 사업시행으로 인하여 건물이 철거되는 이주대상자를 위한 이주대책을 수립하면서 훈령의 형식으로 ‘B시 이주민지원규정’을 마련하였다. \par
위 지원규정에서는 ① 이주대책대상자 선정과 관련하여, 「공익사업을 위한 토지 등의 취득 및 보상에 관한 법률」 및 그 시행령이 정하고 있는 이주대책대상자 요건 외에 ‘전세대원이 사업구역 내 주택 외 무주택’이라는 요건을 추가적으로 규정하는 한편, ② B시의 이주택지 지급 대상에 관하여, 과거 건축물양성화기준일 이전 건물의 거주자의 경우 소지가(조성되지 아니한 상태에서의 토지가격) 분양대상자로, 기준일 이후 건물의 거주자의 경우 일반우선 분양대상자로 구분하고 있는 바, 소지가 분양대상자의 경우 1세대당 상업용지 3평을 일반분양가로 추가 분양하도록 하고, 일반우선 분양대상자의 경우 1세대 1필지 이주택지를 일반분양가로 우선분양할 수 있도록 하고 있다. \par
B시의 시장은 이주대책을 실시하면서 이 사건 개발사업 구역 내에 거주하는 甲과 乙에 대하여, 甲은 공익사업을 위한 토지 등의 취득 및 보상에 관한 법령이 정한 이주대책대상자에 해당됨에도 위 ①에서 정하는 요건을 이유로 이주대책대상자에서 배제하는 부적격 통보를 하였고, 소지가 분양대상자로 신청한 乙에 대해서는 위 지원규정을 적용하여 소지가 분양대상이 아닌 일반우선 분양대상자로 선정하고, 이를 공고하였다. 다음 물음에 답하시오. (30점) \par
\medskip\noindent 1) 甲은 ‘B시 이주민지원규정’에서 정한 추가적 요건을 이유로 자신을 이주대책대상자에서 배제한 것은 위법하다고 주장한다. 甲의 주장이 타당한지에 관하여 설명하시오. (15점) \par
\medskip\noindent 2) 乙은 자신을 소지가 분양대상자가 아닌 일반우선 분양대상자로 선정한 것은 위법하다고 보아 이를 소송으로 다투려고 한다. 乙이 제기하여야 하는 소송의 형식을 설명하시오. (15점) \par
\end{longquestion}
